% Pandoc/Quarto template for den-paper extension
\documentclass[
  manuscript=article,
  layout=publish,
  year=2026,
  volume=1,
]{den}

% DOI

% Conference (for proceedings/poster)

% Dates
\published{20 May 2026}

% Editor and reviewers

% Strip abstract fields from bibliography entries (prevents LuaLaTeX parsing issues)

% Title
\title{User Guide for the DEN Paper Template}

% Authors and affiliations
\author{First Author \orcid{0000-0000-0000-0000}}
\affiliation{Institution-1, City, Country}
\email{correspondence@email.domain}
\author{Krisna Gupta \orcid{0000-0001-8695-8514}}
\affiliation{Dewan Ekonomi Nasional}
\author{Third Author}
\affiliation{Institution-3, City, Country}

% Keywords
\keywords{keyword; keyword-two; keyword number three}

% JEL classification codes
\jel{C22, O47}

% Abbreviations

% Pandoc/Quarto compatibility macros
\providecommand{\tightlist}{%
  \setlength{\itemsep}{0pt}\setlength{\parskip}{0pt}}
\providecommand{\citeproc}[2]{#1}
\providecommand{\citeproctext}{}

% Support for graphics
\RequirePackage{graphicx}

% Support for ding symbols (checkmarks, etc.)
\RequirePackage{pifont}

% Support for table column width calculations (required by Quarto-generated tables)
\RequirePackage{calc}

% Pandoc image handling - required for newer Pandoc/Quarto versions
% Constrains images to text width to prevent overflow
\makeatletter
\providecommand{\pandocbounded}[1]{%
  \sbox0{#1}%
  \ifdim\wd0>\linewidth
    \resizebox{\linewidth}{!}{#1}%
  \else
    #1%
  \fi
}
\makeatother

% CSL references support

% Header includes
\makeatletter
\@ifpackageloaded{caption}{}{\usepackage{caption}}
\AtBeginDocument{%
\ifdefined\contentsname
  \renewcommand*\contentsname{Daftar Isi}
\else
  \newcommand\contentsname{Daftar Isi}
\fi
\ifdefined\listfigurename
  \renewcommand*\listfigurename{Daftar Gambar}
\else
  \newcommand\listfigurename{Daftar Gambar}
\fi
\ifdefined\listtablename
  \renewcommand*\listtablename{Daftar Tabel}
\else
  \newcommand\listtablename{Daftar Tabel}
\fi
\ifdefined\figurename
  \renewcommand*\figurename{Gambar}
\else
  \newcommand\figurename{Gambar}
\fi
\ifdefined\tablename
  \renewcommand*\tablename{Tabel}
\else
  \newcommand\tablename{Tabel}
\fi
}
\@ifpackageloaded{float}{}{\usepackage{float}}
\floatstyle{ruled}
\@ifundefined{c@chapter}{\newfloat{codelisting}{h}{lop}}{\newfloat{codelisting}{h}{lop}[chapter]}
\floatname{codelisting}{Daftar}
\newcommand*\listoflistings{\listof{codelisting}{Daftar Daftar}}
\makeatother
\makeatletter
\makeatother
\makeatletter
\@ifpackageloaded{caption}{}{\usepackage{caption}}
\@ifpackageloaded{subcaption}{}{\usepackage{subcaption}}
\makeatother

% Keep internal links clickable without PDF border underlines.
\hypersetup{pdfborder={0 0 0}}

% Source note environment for figure attributions (right-aligned, compact)
\newenvironment{source}{%
  \vspace{-1em}%
  \begin{flushright}\scriptsize\itshape
}{%
  \end{flushright}%
}
\newcommand{\sumber}[1]{%
  \vspace{-1em}%
  \begin{flushright}\scriptsize\itshape
  Sumber: #1
  \end{flushright}%
}

% Allow LaTeX more flexibility in line breaking to prevent overflow
% (needed for Indonesian/non-English text on narrow B5 paper)
\tolerance=1000
\emergencystretch=1.5em

\begin{document}

\begin{abstract}
An abstract summarizes, in one paragraph of 300 words or fewer, the
major aspects of the entire paper. It should include: 1) the overall
purpose of the study and the research problem you investigated; 2) the
basic design of your research approach; 3) major findings as a result of
your analysis; and 4) a brief summary of your interpretations and
conclusions.
\end{abstract}

\section{Pendahuluan}\label{pendahuluan}

Trade liberalization has long been promoted as a pathway to economic
growth. The classical theories of Ricardo and Heckscher-Ohlin predict
that reducing tariff barriers allows countries to specialize in their
comparative advantage, lowering consumer prices and raising aggregate
output. International institutions including the World Bank, the IMF,
and the WTO have consistently advocated trade openness as a pillar of
development strategy, and Indonesia has followed this counsel through
successive rounds of tariff reduction. From the WTO accession in 1995 to
AFTA, ACFTA, RCEP, and a series of bilateral agreements, Indonesia's
simple average applied tariff has fallen from approximately 15 percent
in the early 1990s to roughly 2 percent by 2023 (WITS, 2024).

Manufacturing needs particular attention in this context because it is
the single largest contributor to GDP growth. In 2024, manufacturing
industries contributed 1.07 percentage points to Indonesia's 5.11
percent GDP growth, exceeding the contributions of agriculture, trade,
mining, and services (BPS, 2024).

Within the manufacturing sector, the question of who benefits from trade
liberalization takes on added urgency. Melitz (2003) demonstrates that
trade liberalization favors large, productive, export-oriented firms,
which gain access to cheaper inputs and larger markets. Small,
domestic-market firms, by contrast, face the downside of increased
import competition without offsetting export opportunities.

\section{Studi Literatur}\label{studi-literatur}

Trade liberalization and welfare has been the subject of a number of
empirical research. Standard trade models predict aggregate welfare
gains from the removal of trade barriers, where liberalization allows
economies to specialise according to comparative advantage and open
consumers to access of more varied goods (Krugman, Obstfeld and Melitz,
2022). These aggregate predictions, however, say little about how the
gains are distributed, and the empirical literature has increasingly
turned to that question.

Much of the recent evidence shows that the direction of the effect
depends on which tariffs are cut. Kis-Katos and Sparrow (2015) combine
the annual manufacturing survey (Survei Industri or IBS) with Sakernas
for the period 1993 to 2002 and estimate effects at both the provincial
and the district level in Indonesia. Both levels of analysis indicate
that average wages rose in response to reductions in tariffs on
intermediate inputs. The two tariff types also move poverty in opposite
directions, since lower output tariffs are associated with higher
poverty across the poverty measures while lower input tariffs are
associated with poverty reduction.

Amiti and Konings (2007) find that a 10 percentage point fall in input
tariffs raises productivity by about 12 percent for firms that import
their inputs, at least twice the gain from an equivalent cut in output
tariffs. To be precise, importers receive this full benefit while
non-importers receive only about 3 percent, so the gains from input
liberalization are concentrated among a subset of firms rather than
shared evenly across the sector.

Evidence from Brazil points to a further source of heterogeneity, this
time across labour market segments. Paz (2014) finds that the gains
accrued disproportionately to formal employment. A one percentage point
reduction in the tariffs applied by trading partners to Brazilian
exports raised formal-sector wages by roughly 0.32 percent, whereas a
one percentage point reduction in Brazil's own import tariffs lowered
formal wages by about 0.05 percent. Lower partner tariffs also reduced
the probability that a worker was employed informally.

The evidence is not uniformly positive. Delesalle et al.~(2025) find no
significant average effect of the ASEAN-China Free Trade Area on the
probability of holding a wage job in Indonesia. The estimated overall
effect is economically very small or statistically insignificant for
men, for women, and for less-educated workers. It is larger and
statistically significant only for educated workers, for whom the
average increase is 1.9 percentage points against a baseline probability
of 17.4 percent in 2000. Since this group accounts for around a quarter
of the workforce, the distributional incidence of the agreement was
regressive. Topalova (2007) reports a sharper result for India, where a
district exposed to the average tariff reduction recorded a poverty
headcount roughly two percentage points higher and a poverty gap roughly
0.6 percentage points higher than a district whose tariffs were
unchanged.

Whether such adverse effects materialise appears to depend on the
institutional setting. Goldberg and Pavcnik (2003) find that Colombian
tariff reductions raised informal employment by roughly 0.1 percentage
points per percentage point of tariff cut, but only before Law 50 of
1990 reduced firing costs. Post-reform, the interaction between tariffs
and informal employment is consistently positive, indicating that the
same tariff decline leads to a smaller increase in informality
afterwards, and possibly a decrease. These results indicate that trade
policy moves informality only when formal employment is costly to
adjust.

Goldberg and Pavcnik (2007) find that globalization increases the wage
skill premium in developing countries at the economy-wide level, within
industries, and within firms. Amiti and Cameron (2012), however, report
the opposite for Indonesian manufacturing. Using firm-level data with
firm and year by island fixed effects, they find that a 10 percentage
point reduction in input tariffs lowered the within-firm skill premium
by about 1.4 to 2.6 percent, while output tariffs had no significant
effect once input tariffs were controlled for. The effect is
concentrated among firms that actually import intermediates.

In analysing the effects of trade liberalization on poverty and labour
market outcomes, Kis-Katos and Sparrow (2015) construct district-level
tariff exposure by combining nationally determined sectoral import
tariffs with each district's initial economic structure prior to trade
liberalization. This approach follows a \textbf{shift-share (Bartik)
framework}, in which national tariff changes are weighted by
predetermined district-level employment or industrial output shares.

For the poverty analysis, district output tariff exposure is measured
using the initial employment distribution across sectors. Specifically,
the tariff exposure of district \(k\) in year \(t\) is defined as

\[
Tariff_{kt}
=
\sum_{s=1}^{S}
\left(
\frac{Q_{sk,1990}}{Q_{k,1990}}
\times
Tariff_{st}
\right),
\]

where:

\begin{itemize}
\tightlist
\item
  \(Q_{sk,1990}\) denotes employment in output sector \(s\) in district
  \(k\) in 1990.
\item
  \(Q_{k,1990}\) is the total labour force in district \(k\) in 1990.
\item
  \(Tariff_{st}\) is the national import tariff applied to sector \(s\)
  in year \(t\).
\item
  \(S\) denotes the total number of tradable sectors.
\end{itemize}

For the firm-level analysis, Kis-Katos and Sparrow (2015) alternatively
construct an output tariff measure using the initial industrial
structure of each district. In this case, sectoral tariffs are weighted
by each sector's share in the district's total industrial output in
1993, \(\left(\frac{Q_{sk,1993}}{Q_{k,1993}}\right)\), rather than by
employment shares.

To capture the effects of tariffs on intermediate inputs, the authors
also construct an input tariff measure. The input tariff for each
industry is first calculated as the weighted average tariff on all
imported intermediate inputs, where the weights are given by each
input's share in the industry's total input use. These industry-specific
input tariffs are then aggregated to the district level using the
district's initial employment shares:

\[
InputTariff_{kt}
=
\sum_{s=1}^{S}
\left(
\frac{Q_{sk,1990}}{Q_{k,1990}}
\times
\sum_{j=1}^{J}
\left(
\frac{M_{js,1990}}{M_{s,1990}}
\times
Tariff_{jt}
\right)
\right).
\]

where:

\begin{itemize}
\tightlist
\item
  \(M_{js,1990}\) denotes the value of intermediate input \(j\) used by
  industry \(s\) in 1990.
\item
  \(M_{s,1990}\) is the total value of intermediate inputs used by
  industry \(s\) in 1990.
\item
  \(\frac{M_{js,1990}}{M_{s,1990}}\) represents the share of input \(j\)
  in the total intermediate inputs of industry \(s\).
\item
  \(Tariff_{jt}\) denotes the national tariff on intermediate input
  \(j\) in year \(t\).
\item
  \(J\) represents the total number of input sectors.
\end{itemize}

Consistent with Kovak (2013), Kis-Katos and Sparrow (2015) exclude
non-tradable sectors from the construction of the output tariff measure
because these sectors are not directly exposed to international trade.

The effects of trade liberalization on poverty and labour market
outcomes are then estimated using the following first-difference
specification:

\[
\Delta y_{kt}
=
\alpha
+
\beta_1
\Delta OutputTariff_{kt}
+
\beta_2
\Delta InputTariff_{kt}
+
\Delta X_{kt}^{\prime}\gamma
+
I_k^{\prime}\theta
+
\lambda_{rt}
+
\Delta\varepsilon_{kt},
\]

where:

\begin{itemize}
\tightlist
\item
  \(\Delta y_{kt}\) denotes the change in the district-level outcome
  variable, including measures of poverty, employment, and wages.
\item
  \(\Delta OutputTariff_{kt}\) represents the change in district-level
  output tariff exposure.
\item
  \(\Delta InputTariff_{kt}\) represents the change in district-level
  input tariff exposure.
\item
  \(\Delta X_{kt}\) is a vector of time-varying control variables,
  including the share of the rural population, the share of the
  working-age population (16--60 years), adult literacy rates, and
  minimum wages.
\item
  \(I_k\) is a vector of initial district characteristics, including the
  1990 labour shares used as tariff weights, aggregated to one-digit
  sectors, the 1990 rural population share, and, in some specifications,
  the initial level of the dependent variable.
\item
  \(\lambda_{rt}\) denotes island-by-year fixed effects.
\item
  \(\Delta\varepsilon_{kt}\) is the error term.
\end{itemize}

The first-difference specification removes time-invariant district
characteristics and exploits within-district variation in tariff
exposure over time. By combining predetermined employment shares with
nationally determined tariff changes, the Bartik framework identifies
heterogeneous district exposure to trade liberalization while reducing
concerns regarding reverse causality.

Unlike Kis-Katos and Sparrow (2015), who estimate the model using a
first-difference OLS specification, Delesalle et al.~(2025) employ a
Bartik instrumental variable (IV)strategy to address the potential
endogeneity of tariff reductions. District-level exposure to trade
liberalization depends on both the pre-liberalization industrial
structure and industry-level tariff changes. Rather than examining
multilateral trade liberalization, the authors focus exclusively on the
bilateral tariff reductions between Indonesia and China under the
ASEAN--China Free Trade Agreement (ACFTA), using pre-ACFTA industry
shares to construct the exposure measures. Specifically, they consider
three channels through which trade liberalization may affect local labor
markets: output import tariff cuts, input import tariff cuts, and export
tariff cuts.

The output import tariff measure is constructed as

\[
\Delta OutputImportTariff_{d,2000-2010}
=
\sum_{s}
\frac{L_{s,d,2000}}{L_{d,2000}}
\cdot
\Delta Tariff^{Indo}_{s,2000-2010}
\]

where:

\begin{itemize}
\tightlist
\item
  \(L_{s,d,2000}\) is the number of working-age individuals employed in
  industry \(s\) in district \(d\) in 2000.
\item
  \(L_{d,2000}\) is the total working-age population employed in
  district \(d\) in 2000.
\item
  \(\Delta Tariff^{Indo}_{s,2000-2010}\) denotes the change in
  Indonesia's import tariff on Chinese products in industry \(s\)
  between 2000 and 2010.
\end{itemize}

To capture the effect of imported intermediate inputs, the authors
construct the input import tariff measure:

\[
\Delta InputImportTariff_{d,2000-2010}
=
\sum_{s}
\frac{L_{s,d,2000}}{L_{d,2000}}
\left(
\sum_{j}
\frac{M_{j,s,2000}}{M_{s,2000}}
\cdot
\Delta Tariff^{Indo}_{j,2000-2010}
\right)
\]

where:

\begin{itemize}
\tightlist
\item
  \(M_{j,s,2000}\) denotes the value of intermediate input \(j\) used by
  industry \(s\).
\item
  \(M_{s,2000}\) is the total intermediate inputs used by industry
  \(s\).
\item
  \(\frac{M_{j,s,2000}}{M_{s,2000}}\) represents the cost share of input
  \(j\) in industry \(s\).
\item
  The inner summation calculates the industry-level input tariff, while
  the outer summation aggregates these tariffs to the district level
  using initial employment shares.
\end{itemize}

The export tariff measure follows the same specification as the output
import tariff but replaces Indonesian import tariffs with Chinese
tariffs imposed on Indonesian exports:

\[
\Delta ExportTariff_{d,2000-2010}
=
\sum_{s}
\frac{L_{s,d,2000}}{L_{d,2000}}
\cdot
\Delta Tariff^{China}_{s,2000-2010}
\]

where \(\Delta Tariff^{China}_{s,2000-2010}\) denotes the change in
Chinese import tariffs on Indonesian exports in industry \(s\).

The relationship between tariff exposure and labor market outcomes is
estimated using the following empirical specification:

\[
\Delta y_{d,2000-2010}
=
\alpha
+
\beta_1
\Delta OutputImportTariff_{d}
+
\beta_2
\Delta InputImportTariff_{d}
+
\beta_3
\Delta ExportTariff_{d}
+
\gamma X_{d,2000}
+
\varepsilon_d
\]

where:

\begin{itemize}
\tightlist
\item
  \(\Delta y_{d,2000-2010}\) denotes the change in the district-level
  labor market outcome between 2000 and 2010.
\item
  \(\Delta OutputImportTariff_{d}\) is the district-level output import
  tariff exposure.
\item
  \(\Delta InputImportTariff_{d}\) is the district-level input import
  tariff exposure.
\item
  \(\Delta ExportTariff_{d}\) is the district-level export tariff
  exposure.
\item
  \(X_{d,2000}\) is a vector of pre-liberalization district
  characteristics, including industrial composition, education,
  demographic variables, and island fixed effects.
\item
  \(\varepsilon_d\) is the error term.
\end{itemize}

To address the potential endogeneity of tariff reductions, Delesalle et
al.~(2025) instrument district-level tariff changes using the initial
average tariff level in 2000. The identification strategy exploits the
design of the ACFTA, under which sectors with higher pre-liberalization
tariffs experienced systematically larger tariff reductions between 2000
and 2010. Consequently, the initial tariff level provides a strong
predictor of subsequent tariff changes while remaining plausibly
exogenous to later labor market outcomes after conditioning on
pre-treatment characteristics.

Unlike Kis-Katos and Sparrow (2015), who rely on a first-difference OLS
specification, Delesalle et al.~(2025) estimate the causal effects using
a Bartik IV strategy, thereby strengthening causal identification by
mitigating potential bias arising from the non-random allocation of
tariff reductions across industries.

\section{Metodologi}\label{metodologi}

\subsection{Data}\label{data}

\subsubsection{Data and Unit of
Analysis}\label{data-and-unit-of-analysis}

This study utilizes the Survei Industri Mikro dan Kecil (IMK)
administered by Badan Pusat Statistik (BPS) for the period 2010 to 2024.
The IMK covers manufacturing establishments with 1 to 19 workers. This
distinguishes it from the Survei Industri Besar dan Sedang (IBS), which
enumerates establishments with 20 or more workers).

The period 2010 to 2024 is selected to capture the effects of continued
tariff reduction on micro and small manufacturing. xxxx

An important measurement constraint should be noted at the outset. The
IMK records workers only as paid or unpaid, and this classification does
not correspond to informality as defined in the comparative literature.
Goldberg and Pavcnik (2003) define informality through compliance with
labour market regulation, a definition the IMK cannot support because it
collects no information on contracts, social security enrolment, or
regulatory registration.

Tariff data are drawn from the World Integrated Trade Solution (WITS)
database maintained by the World Bank, using UNCTAD-TRAINS applied rates

The unit of analysis is the IMK firms at the province-year level.

\subsubsection{Distribution Sample}\label{distribution-sample}

The provincial distribution of IMK firms during the 2018--2022 period
reveals substantial variation across Indonesia. West Java recorded the
highest number of IMK firms in the sample, with 65 firms, indicating
that the province has the greatest concentration of small manufacturing
establishments among all provinces included in the dataset. \textbar{}
\textbf{Province} \textbar{} \textbf{Number of IMK Firms} \textbar{}
\textbar:----------\textbar--------------------:\textbar{} \textbar{}
Aceh \textbar{} 7 \textbar{} \textbar{} North Sumatra \textbar{} 16
\textbar{} \textbar{} West Sumatra \textbar{} 4 \textbar{} \textbar{}
Riau \textbar{} 29 \textbar{} \textbar{} Jambi \textbar{} 1 \textbar{}
\textbar{} Lampung \textbar{} 13 \textbar{} \textbar{} Riau Islands
\textbar{} 5 \textbar{} \textbar{} DKI Jakarta \textbar{} 23 \textbar{}
\textbar{} West Java \textbar{} 65 \textbar{} \textbar{} Central Java
\textbar{} 35 \textbar{} \textbar{} DI Yogyakarta \textbar{} 18
\textbar{} \textbar{} East Java \textbar{} 61 \textbar{} \textbar{}
Banten \textbar{} 8 \textbar{} \textbar{} Bali \textbar{} 9 \textbar{}
\textbar{} West Nusa Tenggara \textbar{} 4 \textbar{} \textbar{} West
Kalimantan \textbar{} 2 \textbar{} \textbar{} Central Kalimantan
\textbar{} 4 \textbar{} \textbar{} East Kalimantan \textbar{} 10
\textbar{} \textbar{} North Sulawesi \textbar{} 10 \textbar{} \textbar{}
Central Sulawesi \textbar{} 1 \textbar{} \textbar{} South Sulawesi
\textbar{} 1 \textbar{} \textbar{} West Papua \textbar{} 8 \textbar{}
\textbar{} Papua \textbar{} 9 \textbar{} \#\# Metode

\section{Hasil}\label{hasil}

\section{Kesimpulan}\label{kesimpulan}

\section{Referensi}\label{referensi}

\subsection{References}\label{references}

Add your bibliography entries to \texttt{ref.bib} and cite them using
standard Quarto citation syntax. Here is an example citation on diamond
open access (\textbf{fuchs2013diamond?}). According to
(\textbf{fuchs2013diamond?}), diamond open access model is amazing.

\subsection{Footnotes}\label{footnotes}

Use footnotes sparingly. They should be used for additional information
that is not essential to the main text.\footnote{This is an example
  footnote.}

\subsection{Tables}\label{tables}

Use standard markdown tables or Quarto's table syntax.
Tabel~\ref{tbl-example} shows an example.

\begin{longtable}[]{@{}lll@{}}
\caption{Example table}\label{tbl-example}\tabularnewline
\toprule\noalign{}
\textbf{Parameter} & \textbf{Notation} & \textbf{Remarks} \\
\midrule\noalign{}
\endfirsthead
\toprule\noalign{}
\textbf{Parameter} & \textbf{Notation} & \textbf{Remarks} \\
\midrule\noalign{}
\endhead
\bottomrule\noalign{}
\endlastfoot
name & - & engine common identifier \\
manufacture & - & name of the manufacturer \\
bpr & \(\lambda\) & bypass ratio \\
pr & - & pressure ratio \\
thrust & \(T_0\) & maximum static thrust \\
\end{longtable}

\subsection{Figures}\label{figures}

Use standard Quarto figure syntax with
\texttt{!{[}caption{]}(path)\{\#fig-id\ width=X\%\}}. Place images in a
\texttt{fig/} folder.

\subsection{Equations}\label{equations}

Use standard LaTeX equation syntax. Persamaan~\ref{eq-cauchy} shows an
example equation.

\begin{equation}\protect\phantomsection\label{eq-cauchy}{
\rho\frac{\mathrm{D} \mathbf{u}}{\mathrm{D} t} = - \nabla p + \nabla \cdot \boldsymbol \tau + \rho\,\mathbf{g}
}\end{equation}

\section{Sections}\label{sections}

Organize your paper using standard markdown headings.

Some standard sections are:

\begin{itemize}
\tightlist
\item
  Introduction
\item
  Methods
\item
  Results
\item
  Discussion
\item
  Conclusion
\end{itemize}

\section*{Acknowledgement}\label{acknowledgement}
\addcontentsline{toc}{section}{Acknowledgement}

Include your acknowledgements in this section.

\section*{Funding Statement}\label{funding-statement}
\addcontentsline{toc}{section}{Funding Statement}

When applicable, please specify the funding information for this
research.

\section*{Reproducibility Statement}\label{reproducibility-statement}
\addcontentsline{toc}{section}{Reproducibility Statement}

\textbf{Mandatory section!}

\subsection{Information on how to reproduce this research, including
access to: 1) source code related to the research, 2) source code for
the figures, 3) source code/data for the tables when
applicable.}\label{information-on-how-to-reproduce-this-research-including-access-to-1-source-code-related-to-the-research-2-source-code-for-the-figures-3-source-codedata-for-the-tables-when-applicable.}

\section{Pendahuluan}\label{pendahuluan-1}

\section{Studi Literatur}\label{studi-literatur-1}

xxxx

\section{Metodologi}\label{metodologi-1}

\subsection{Data}\label{data-1}

\subsection{Metode}\label{metode}

\section{Analisis}\label{analisis}

\end{document}
